\documentclass[11pt]{article}
\usepackage[margin=1in]{geometry}
\usepackage{amsmath,amssymb,amsthm,mathtools}
\usepackage{graphicx}
\usepackage[hidelinks]{hyperref}
\usepackage{microtype}

\newtheorem{theorem}{Theorem}
\newtheorem{lemma}[theorem]{Lemma}
\newtheorem{remark}[theorem]{Remark}
\newcommand{\HH}{\mathcal H}
\newcommand{\B}{\mathcal B}
\newcommand{\Lone}{\mathcal L^1}
\newcommand{\U}{\mathcal U}
\newcommand{\Tr}{\operatorname{Tr}}
\newcommand{\essspec}{\sigma_{\mathrm{ess}}}
\newcommand{\clos}[1]{\overline{#1}}
\newcommand{\WC}[2][]{W_{#1C}(#2)}

\title{Convexity of the ordinary infinite-dimensional
\texorpdfstring{$C$}{C}-numerical range for positive
\texorpdfstring{$C$}{C} with finite kernel}
\author{Candidate partial result for arXiv:2009.01300\\
Lane 12, agent \texttt{agent\_lane\_12}, model GPT5.6}
\date{22 July 2026}

\begin{document}
\maketitle

\begin{center}
\fbox{\parbox{0.92\linewidth}{
\textbf{Status: partial result, likely valid; human verification requested.}
The proof gives a complete positive answer when $C\geq0$ is trace class with
$\dim\ker C<\infty$, for every bounded $A$.  The general selfadjoint problem,
and the positive infinite-rank/infinite-kernel case for arbitrary $A$, remain
open in this packet.}}
\end{center}

\section{Source question and result}

Let $\HH$ be an infinite-dimensional separable complex Hilbert space.  For
$C\in\Lone(\HH)$ and $A\in\B(\HH)$, the ordinary $C$-numerical range is
\[
  W_C(A):=\{\Tr(UCU^*A):U\in\U(\HH)\}.
\]
Dirr and vom Ende asked whether $W_C(A)$ is convex whenever $C$ is
selfadjoint~\cite[Open Problem (b)]{DvE}.  Loreaux and Patnaik restated the
question and proved convexity for their orbit-closed modification, together
with ordinary convexity under special hypotheses~\cite{LP}.  The question is
recorded in the abstract of arXiv:2009.01300, reproduced below.

\begin{figure}[ht]
  \centering
  \includegraphics[width=\linewidth]{figures/open_problem_crop.png}
  \caption{Full-width crop of PDF page 1 of arXiv:2009.01300.  The abstract
  says that Dirr--vom Ende asked whether the ordinary $C$-numerical range is
  convex without taking closure, and describes the paper's result as only a
  partial answer for the ordinary range.}
\end{figure}

\noindent\textbf{Transcription of the question.}
For selfadjoint trace-class $C$ and bounded $A$, is $W_C(A)$ convex without
taking its closure?

Our result settles the finite-kernel positive subcase without any hypothesis
on $A$.

\begin{theorem}[Finite-kernel positive case]\label{thm:main}
Let $C\in\Lone(\HH)$ be positive and suppose $\dim\ker C<\infty$.  Then
$W_C(A)$ is convex for every $A\in\B(\HH)$.
\end{theorem}

Because $\HH$ is infinite-dimensional, such a $C$ has infinite rank.  The
finite-rank case is already known from Westwick's theorem~\cite{Westwick}.

\section{Proof intuition}

There are two different geometric regimes.  First, every point in the planar
interior of $\clos{W_C(A)}$ is already in $W_C(A)$.  To see this, surround the
point by three orbit values, approximate the corresponding three orbit
operators so that they share one common infinite-dimensional tail, and vary
only a common finite-dimensional head.  The finite-dimensional
$C$-numerical range of that head is convex, so it fills the triangle exactly.

Only boundary segments remain.  After rotating a supporting line, points on
such a segment are ordinary maximizers of $X\mapsto\Tr(XH)$ for a selfadjoint
$H$.  The equality analysis of Loreaux--Patnaik forces every maximizer to put
the same successive eigenvalue blocks of $C$ into the spectral blocks of $H$
above the essential spectral maximum.  It vanishes below that maximum.  The
finite-kernel assumption fixes the kernel multiplicity of the sole residual
block at the threshold.  Hence any two ordinary maximizers are conjugate by a
unitary commuting with $H$.  A logarithm path in that commutant stays on the
supporting face, and the intermediate value theorem fills the segment.

\section{Interior points}

The first lemma is stated for arbitrary selfadjoint trace-class $C$.

\begin{lemma}[Common-tail approximation]\label{lem:tail}
Let $C=C^*\in\Lone(\HH)$ and let $X_j=U_jCU_j^*$, $1\leq j\leq r$, be a
finite family in its unitary orbit.  For every $\varepsilon>0$ there are a
finite-dimensional subspace $M\subset\HH$, a single trace-class operator
$D$ on $M^\perp$, and orbit elements
\[
  Y_j=B_j\oplus D \qquad (\HH=M\oplus M^\perp)
\]
such that $\|X_j-Y_j\|_1<\varepsilon$ for all $j$, and all $B_j$ belong to
one finite-dimensional selfadjoint unitary orbit in $\B(M)$.
\end{lemma}

\begin{proof}
Choose a finite-rank spectral truncation $C_N$ of $C$ such that
$2\|C-C_N\|_1<\varepsilon$, and let $E_N$ be its supporting
finite-dimensional spectral subspace.  Choose a finite-dimensional $M$
containing all $U_jE_N$ and put $d=\dim M$.  There is a $C$-invariant
$d$-dimensional subspace $E_d\supset E_N$: extend by further nonzero
eigenvectors, and then by kernel vectors if necessary.  Choose an eigenbasis
$e_1,\ldots,e_d$ of $E_d$ whose initial vectors support $C_N$.

For each $j$, extend the prescribed orthonormal family obtained from those
initial vectors under $U_j$ to a basis $f_{j,1},\ldots,f_{j,d}$ of $M$.
Fix a unitary $R:E_d^\perp\to M^\perp$, and define a unitary $V_j$ by
\[
  V_je_n=f_{j,n}\quad(1\leq n\leq d),
  \qquad V_j|_{E_d^\perp}=R.
\]
Set $Y_j=V_jCV_j^*$.  The $C_N$ parts of $X_j$ and $Y_j$ coincide, whence
\[
 \|X_j-Y_j\|_1
 \leq \|U_j(C-C_N)U_j^*\|_1+\|V_j(C-C_N)V_j^*\|_1
 =2\|C-C_N\|_1<\varepsilon.
\]
Relative to $M\oplus M^\perp$, the head $B_j$ has the eigenvalue list of
$C|_{E_d}$, independent of $j$, and the tail is the common operator
$D=R(C|_{E_d^\perp})R^*$.
\end{proof}

\begin{lemma}[Interior filling]\label{lem:interior}
If $C=C^*\in\Lone(\HH)$ and $A\in\B(\HH)$, then
\[
  \operatorname{int}\clos{W_C(A)}\subseteq W_C(A),
\]
where the interior is the ordinary planar interior.
\end{lemma}

\begin{proof}
Loreaux--Patnaik prove that the orbit-closed range is convex
and that $W_C(A)$ is dense in it~\cite[Theorem 3.4 and Corollary 4.6]{LP}.
Thus $K:=\clos{W_C(A)}$ is convex.  Let $z\in\operatorname{int}K$.
Choose three values $z_j=\Tr(X_jA)$ in $W_C(A)$ whose triangle contains $z$
in its strict interior.  This is possible by first taking a small triangle
around $z$ inside $K$, then perturbing its vertices into the dense subset
$W_C(A)$.

Apply Lemma~\ref{lem:tail} with error small enough that
$y_j:=\Tr(Y_jA)$ still have a triangle containing $z$; indeed
$|y_j-z_j|\leq\|Y_j-X_j\|_1\|A\|$.  Let $A_M=P_MA|_M$.  Keeping the common
tail $D$ fixed and varying the finite head through its unitary orbit gives
exactly the subset
\[
  \Tr(DP_{M^\perp}A|_{M^\perp})+W_{C_M}(A_M)
  \subseteq W_C(A),                                      \tag{3.1}
\]
where $C_M=C|_{E_d}$ transported to $M$.  Off-diagonal blocks of $A$ make no
contribution to the trace.  The set in (3.1) is convex by the
finite-dimensional $C$-numerical-range theorem for selfadjoint $C_M$
\cite{Westwick,Poon}.  It contains all three $y_j$, hence their triangle and
therefore $z$.
\end{proof}

\section{Supporting faces}

We now assume $C\geq0$ and put $k=\dim\ker C<\infty$.  Denote by
$\mathcal O(C)$ the trace-norm closure of the unitary orbit.

\begin{lemma}[One commutant orbit of ordinary maximizers]
\label{lem:maximizers}
Let $H=H^*\in\B(\HH)$.  If $X_1,X_2\in\U(C)$ both maximize
$X\mapsto\Tr(XH)$ over $\U(C)$, then there is a unitary $V$ commuting with
$H$ such that $X_2=VX_1V^*$.
\end{lemma}

\begin{proof}
Let $m=\max\essspec(H)$ and $P=\chi_{[m,\infty)}(H)$.  Density of the
ordinary range in the orbit-closed range makes the two suprema equal.
Therefore the maximizer-support theorem
\cite[Theorem 5.7]{LP} applies and gives
\[
  X_j=PX_jP.                                               \tag{4.1}
\]
Since $X_j$ is unitarily equivalent to $C$, its kernel has dimension $k$.
Thus $P^\perp\HH\subseteq\ker X_j$, so
\[
  q:=\dim P^\perp\leq k<\infty.                           \tag{4.2}
\]

List the distinct eigenvalues of $H$ strictly above $m$ in decreasing order,
and let $P_\ell$ be the associated finite-dimensional spectral projections.
The block-decomposition equality theorem
\cite[Proposition 5.12]{LP} says that each $X_j$ commutes with every
$P_\ell$, and its compression to $P_\ell\HH$ is unitarily equivalent to the
same successive finite block of the decreasing positive eigenvalues of $C$.

Suppose first that there are infinitely many such blocks.  Their cumulative
dimensions tend to infinity, so the prescribed blocks exhaust every nonzero
eigenvalue of $C$.  Their traces therefore sum to $\Tr C=\Tr X_j$.
Positivity forces $X_j$ to vanish on the orthogonal complement of their
direct sum.  The two compressions are unitarily conjugate on each block and
both are zero on the complement.

Suppose instead that there are only finitely many blocks strictly above $m$.
Proposition 5.12 places the remaining nonzero eigenvalue tail of $C$ on the
threshold space $P_0\HH$, where $P_0=\chi_{\{m\}}(H)$, while (4.1) makes
$X_j$ zero on $P^\perp\HH$.  The blocks strictly above $m$ contain only
positive eigenvalues and have no kernel.  Hence
\[
 \dim\ker(X_j|_{P_0\HH})=k-q,                              \tag{4.3}
\]
the same finite number for $j=1,2$.  The two residual positive compact
operators have the same nonzero eigenvalues with multiplicity and the same
kernel dimension, so the spectral theorem makes them unitarily equivalent.
If $P_0=0$, the existence of $X_j$ simply forces the residual tail to be
empty, and the same conclusion holds.

Choose a conjugating unitary on every nontrivial spectral block, take their
orthogonal direct sum, and use the identity where both $X_j$ vanish.  The
resulting unitary $V$ conjugates $X_1$ to $X_2$ and commutes with $H$, since
each nontrivial summand lies within an eigenspace of $H$.
\end{proof}

\begin{lemma}[Occupied supporting faces are filled]\label{lem:face}
For $A\in\B(\HH)$, the intersection of $W_C(A)$ with any supporting line of
$K=\clos{W_C(A)}$ is convex.
\end{lemma}

\begin{proof}
After multiplying $A$ by a scalar of modulus one, take the supporting line to
be vertical.  Write $H=\operatorname{Re}A$ and $G=\operatorname{Im}A$.
Two points of $W_C(A)$ on this line have the form
\[
  \Tr(X_jA)=\alpha+i\beta_j,
\]
where $X_j\in\U(C)$ maximize $\Tr(XH)$.  By
Lemma~\ref{lem:maximizers}, $X_2=VX_1V^*$ for a unitary $V$ commuting with
$H$.

Choose a bounded selfadjoint Borel-functional-calculus logarithm $L$ of $V$.
Because $V$ commutes with $H$, so does $L$.  For
$V_t=e^{itL}$ and $X_t=V_tX_1V_t^*$, every $X_t$ lies in $\U(C)$ and
\[
  \Tr(X_tH)=\Tr(X_1H)=\alpha.
\]
The real-valued function $t\mapsto\Tr(X_tG)$ is continuous and has endpoint
values $\beta_1,\beta_2$.  It attains every intermediate value, so the entire
segment between the two original points lies in $W_C(A)$.
\end{proof}

\section{Proof of Theorem~\ref{thm:main}}

Put $S=W_C(A)$ and $K=\clos S$.  The set $K$ is convex by
\cite[Theorem 3.4 and Corollary 4.6]{LP}.  The unitary group is norm
path-connected: if $U$ is unitary, a bounded selfadjoint logarithm produces
the path $e^{it\log U}$.  Hence $S$, a continuous image of the unitary group,
is path-connected.

If $K$ has affine dimension at most one, $S$ is a connected subset of a line
(or a point), hence an interval and therefore convex.  Suppose $K$ has
nonempty planar interior.  Take $x,y\in S$ and $z=(1-t)x+ty$ with $0<t<1$.
If $z\in\operatorname{int}K$, Lemma~\ref{lem:interior} gives $z\in S$.
If $z\in\partial K$, choose a supporting real-linear functional $\phi$ at
$z$.  Since $\phi\leq\phi(z)$ on $K$ and
\[
  \phi(z)=(1-t)\phi(x)+t\phi(y),
\]
both inequalities must be equalities.  Thus $x$ and $y$ lie on the same
supporting line, and Lemma~\ref{lem:face} gives $z\in S$.  Every segment with
endpoints in $S$ is contained in $S$, proving the theorem. \qed

\section{Verification, limitations, and novelty}

\begin{itemize}
\item The common-tail construction was checked at the operator level: the
leading spectral truncations agree exactly, while the two discarded tails
give the explicit $2\|C-C_N\|_1$ bound.  The finite-dimensional convex subset
in (3.1) consists of genuine members of the original unitary orbit.
\item The only place where $\dim\ker C<\infty$ is used is (4.2)--(4.3).
It prevents kernel multiplicity from moving invisibly between the threshold
space and its complement.
\item The proof does not claim closedness.  It fills the planar interior and
each supporting face that already contains ordinary orbit points.
\item When $\dim\ker C=\infty$ and $C$ has infinite rank, cardinal
subtraction does not determine the residual kernel multiplicity in (4.3).
The commutant-orbit lemma is therefore not proved by this argument.  The
general selfadjoint problem also remains open here.
\item A bounded novelty search on 22 July 2026 used arXiv/web queries for
``$C$-numerical range convex'', ``positive trace class'',
``finite-dimensional kernel'', and the Dirr--vom Ende open problem.  It found
arXiv:1712.01023, 1808.06898, 2009.01300, and the orbit-closed follow-up
arXiv:2106.11205, but no later theorem with this finite-kernel conclusion and
arbitrary bounded $A$.  This is evidence, not a definitive literature
certification.
\end{itemize}

\noindent\textbf{Human-review recommendation.}
Verify and retain as a substantial partial theorem.  The main points to audit
are the simultaneous common-tail construction in Lemma~\ref{lem:tail} and the
finite residual-kernel count in (4.3) when Proposition 5.12 of
\cite{LP} is applied to ordinary maximizers.

\begin{thebibliography}{9}
\bibitem{DvE}
G.~Dirr and F.~vom Ende,
\emph{The $C$-numerical range in infinite dimensions},
Linear and Multilinear Algebra 68 (2020), 652--678;
arXiv:1712.01023.

\bibitem{LP}
J.~Loreaux and S.~Patnaik,
\emph{Convexity of the orbit-closed $C$-numerical range and majorization},
arXiv:2009.01300 (2020).

\bibitem{Westwick}
R.~Westwick,
\emph{A theorem on numerical range},
Linear and Multilinear Algebra 2 (1975), 311--315.

\bibitem{Poon}
Y.~T.~Poon,
\emph{Another proof of a result of Westwick},
Linear and Multilinear Algebra 9 (1980), 35--37.
\end{thebibliography}

\end{document}
